\section{The dismantling hierarchy}\label{sec:prelim}
All graphs are finite, simple and undirected. For $S\subseteq V(H)$, write
$H[S]$ for the induced subgraph and $H-v=H[V(H)\setminus\{v\}]$.
The open and closed neighbourhood sets are
\[
 N_H(v)=\{u:uv\in E(H)\},\qquad N_H[v]=N_H(v)\cup\{v\}.
\]
We write $L_H(v)=H[N_H(v)]$ for the \defi{neighbourhood graph}.
A \defi{cone} is a non-empty graph with a universal vertex, called an
\defi{apex}. Thus $K_1$ is a cone and the empty graph is not.

A vertex $u$ is \defi{dominated} by $w\ne u$ when
$N_H[u]\subseteq N_H[w]$. Deleting it is a \defi{fold}, written $u\to w$.
All neighbourhoods in a sequence of folds are taken in the graph then remaining.

\begin{lemma}\label{lem:dom}
A vertex $u$ is dominated if and only if $L_H(u)$ is a cone.
\end{lemma}
\begin{proof}
The inclusion $N_H[u]\subseteq N_H[w]$ says exactly that $w$ is a
neighbour of $u$ adjacent to all its other neighbours.
\end{proof}

The hierarchy of Fieux and Jouve~\cite{FJ2020} replaces the cone test by
a reduction in the preceding class. At level one these are the
$s$-dismantling moves of~\cite{BFJ2010}.
\begin{definition}\label{def:hierarchy}
Let $\D_{-1}$ consist of the non-empty cones. For $j\geq0$, a graph is
in $\D_j$ if it reduces to $K_1$ by successive deletions whose current
neighbourhood graphs belong to $\D_{j-1}$. Such a graph has
\defi{strict level $j$} when it is outside $\D_{j-1}$.
The empty graph belongs to none of these classes.
\end{definition}
The empty reduction puts $K_1$ in every class. Membership is invariant under
isomorphism, since an isomorphism preserves every neighbourhood test.

\begin{lemma}\label{lem:inclusions}
The classes are nested: $\D_{-1}\subseteq\D_0\subseteq\D_1\subseteq\cdots$.
\end{lemma}
\begin{proof}
A cone folds onto its apex. If $\D_{j-1}\subseteq\D_j$, every
neighbourhood in a level-$j$ reduction also passes the level-$(j+1)$ test.
The same sequence proves the next inclusion.
\end{proof}

\begin{lemma}\label{lem:first}
For $|V(H)|>1$ and $j\geq0$,
\begin{equation}\label{eq:first}
 H\in\D_j\quad\Longleftrightarrow\quad
 \text{some }v\in V(H)\text{ satisfies }L_H(v)\in\D_{j-1}
 \text{ and }H-v\in\D_j.
\end{equation}
\end{lemma}
\begin{proof}
Read the first deletion of a reduction for the forward implication.
For the converse, prepend the permitted deletion of $v$ to a reduction of $H-v$.
\end{proof}
In particular, failure of the neighbourhood test at every vertex excludes
membership. Equation~\eqref{eq:first}, together with the empty graph and cone
cases, is also the finite recursion used by the verification programs.
